% latexmk -pdf -pvc blatt11.tex
\documentclass{hott-übung}

\begin{document}

\setcounter{blattnummer}{11}
\setcounter{aufgabennummer}{0}

\blatt

\aufgabe{}
Sei \(A\) ein zusammenhängender \(H\)-Raum mit Multiplikation \(\mu : A \times A \to A\).
Zeige, dass für jedes \(a : A\) die Abbildung \(\mu(a,\_) : A \to A\) eine Äquivalenz ist.

\aufgabe{}
Zeige, dass \(S^0 \ast A \simeq \Sigma A\) gilt.

\aufgabe{Die kleine Hopf-Faserung}
Konstruiere eine \(H\)-Raumstruktur \((S^0,\mu,e)\), sodass \(\mu(x,\_)\) für alle \(x:S^0\) eine Äquivalenz ist. Konstruiere nun analog zur Vorlesung eine Fasersequenz der Form
\[\begin{tikzcd}S^0 \arrow[r] & S^1 \arrow[r] & S^1.\end{tikzcd}\]
Vergleiche zudem die assoziierte Hopf-Faserung mit der \enquote{Orientierungsüberlagerung} \(D : S^1 \to \mU\) aus Beispiel 2.6.5.

\aufgabe{}
Sei \(S\equiv(A \xleftarrow{f} C \xrightarrow{g} B)\) ein Spann.
Der Typ der Kokegel unter \(S\) mit Spitze \(Y\) sei definiert als
\[\Cocone_S(Y) \colonequiv (k : A \to Y) \times (l : B \to Y) \times (K : k\circ f \sim l \circ g).\]
Ein kommutatives Quadrat
\begin{equation}
  \label{square}
  \begin{tikzcd}
  	C & B \\
  	A & X
  	\arrow["g", from=1-1, to=1-2]
  	\arrow["f"', from=1-1, to=2-1]
  	\arrow["H"{description}, draw=none, from=1-1, to=2-2]
  	\arrow["i", from=1-2, to=2-2]
  	\arrow["j"', from=2-1, to=2-2]
  \end{tikzcd}
\end{equation}
heißt \emph{kokartesisch} (auch genannt \emph{Pushout-Quadrat}), wenn die kanonische Abbildung
\[\mathrm{cocone} : (X \to Y) \to \Cocone_S(Y), \quad h \mapsto (h\circ i,h\circ j,h(H))\]
eine Äquivalenz ist.
Zeige, dass ein kommutatives Quadrat wie in \ref{square} genau dann kokartesisch ist, wenn für jedes \(Y : \mU\) das folgende Quadrat kartesisch ist, wobei \(H^*\colonequiv h \mapsto \funext(h(H))\) ist.
\[\begin{tikzcd}
	{Y^X} & {Y^B} \\
	{Y^A} & {Y^C}
	\arrow["{\_\circ j}", from=1-1, to=1-2]
	\arrow["{\_\circ i}"', from=1-1, to=2-1]
	\arrow["{\_\circ g}", from=1-2, to=2-2]
	\arrow["{H^*}"{description}, draw=none, from=2-1, to=1-2]
	\arrow["{\_\circ f}"', from=2-1, to=2-2]
\end{tikzcd}\]

\end{document}
